

\chapter{Compositional differentiation for finite element analysis}




\section{Introduction}

\subsection{Premise}

\begin{itemize}
    \item The direct differentiation method yields sensitivities with respect to parameters that appear natively inside the finite element kernel.
    \item Modern analysis workflows often introduce design variables at a higher level. Typical examples are geometric dimensions, section generators, discretization rules, and load parameterizations.
    \item Implementing general sensitivity by DDM for fiber sections can incur a considerable computational cost
    \item The proposed formulation separates these layers. The kernel differentiates the native mechanics response. The host differentiates the parameterization that constructs native parameters. The full derivative is then recovered by composition.
\end{itemize}


Literature:
\begin{itemize}
    \item \cite{mcdonnell2024geometrically}, \cite{brüls2008sensitivity}, \cite{arnold2007convergence}
    \item \cite{maclaurin2016modeling}
\end{itemize}

% \subsection{Motivation}
% The direct differentiation method (DDM) provides exact sensitivities
% of the discrete equilibrium equations that govern a finite element
% model. 
% Given a residual $\residual(\ModelState;\corey)=\mathbf{0}$
% parameterized by a vector of model inputs~$\corey$, the DDM
% differentiates this equation implicitly with respect to each
% component~$y_i$ and solves the resulting linear system to obtain
% $\partial\ModelState/\partial y_i$ at negligible marginal cost once the
% primal factorization is available.  
% The method is well established in
% computational mechanics and forms the backbone of gradient-based
% design optimization, reliability analysis, and parameter identification
% in structural and solid-mechanics codes.

% Despite its efficiency, the DDM is limited in scope: it can only
% differentiate with respect to quantities that appear explicitly in the
% finite element residual.  In practice, however, the
% parameters of interest to the analyst are often higher-level quantities
% that \emph{determine} these core inputs through an arbitrarily complex pre-processing chain.  
% Cross-section dimensions must be converted to
% fiber discretizations; topology descriptors must be mapped to element
% connectivities and material layouts; stochastic field coefficients
% must be transformed into spatially varying material properties.  Each
% such transformation introduces a link in the chain rule that the DDM alone cannot traverse.

% Automatic differentiation (AD) frameworks such as \Jax{} and PyTorch
% face the complementary limitation.  
% They can differentiate arbitrary
% host-language code with ease, but they are unaware of the implicit
% function defined by the finite element solve.  Treating the analysis
% as an opaque subroutine and differentiating through its internal
% solution procedure is possible in principle \citep{beck1994automatic,gilbert1992automatic,griewank1993derivative} but is neither efficient nor
% numerically robust for large-scale structural models.
% Furthermore, the natural host environments for these frameworks are
% generally not well suited for the efficient implementation of finite
% element procedures.  
% Python, for instance, does not provide sufficient
% control over memory layout and allocation to support the kind of
% low-level optimization that production finite element codes require.
% Even higher-level compiled languages that target scientific computing,
% such as Julia, lack the fine-grained control over memory ownership,
% lifetime, and the separation of compile-time and run-time computation
% that languages like C++ and Fortran treat as first-class abstractions.

% The goal of this chapter is to describe a compositional framework that
% bridges these two worlds.  
% The finite element kernel exposes its
% primal solve and DDM sensitivities as a \emph{differentiable
% primitive} whose forward- and reverse-mode rules are known.  
% The host
% platform contributes automatic differentiation of everything outside
% the kernel.  The chain rule then composes the two layers, yielding
% end-to-end sensitivities with respect to any host-level parameter
% without manual derivation.


% 

Algorithmic differentiation  is a class of computational methods
which leverage the chain rule in order to exactly evaluate\footnote{Within
  machine precision.} the gradient of a mathematical computation which is
implemented as a computer program. 
These methods generally leverage either (1) source code transformation
techniques or (2) operator overloading so that such a gradient
computation can be deduced entirely from the natural semantics of the
programming language being used. 
Much of the early literature which explores applications of \autodiff{} to fields like structural modeling
\autocite{ozaki1995higherorder,martins2001connection} employ the first
approach, where prior to compilation, an pre-processing step is
performed on some source code which generates new code implementing the
derivative of an operation. This approach is generally used with
primitive languages like FORTRAN and C which lack built-in abstractions
like operator overloading and was likely selected for these works
because at the time, such languages were the standard in the field.
Although these compile-time methods allow for powerful optimizations to
be used, they can add significant complexity to a project. 
Languages like C++, Python and Matlab allow user-defined types to \emph{overload}
built-in primitive operators like \texttt{+} and \texttt{*} so that when
they are applied to certain data types they invoke a user-specified
procedure. 
These features allow for \autodiff{} implementations which are
entirely contained within a programming language's standard ecosystem.

\section{Literature}
These methods are often further
categorized into \emph{forward} and \emph{reverse} modes, the later of
which further specializes to the \emph{backpropagation} class of
algorithms which has become ubiquitous in the field of machine learning.
% Hybrid Differentiation in Finite Element Analysis\label{hybrid-differentiation-in-finite-element-analysis}

\subsection{Direct Differentiation in
Computational Mechanics}

The direct differentiation method (DDM) computes sensitivities
by analytically differentiating the finite element governing equations
with respect to model parameters. 
In the earthquake engineering and
reliability community, DDM has been a preferred approach for obtaining
exact response gradients. 
\cite{haukaas2003finite} implemented a comprehensive DDM-based sensitivity module in the \OpenSees{} platform, enabling accurate and efficient computation of
response sensitivities for material, load, and shape parameters.

In such implementations, one derives the partial derivative of each
element's stiffness matrix, internal force, etc., and assembles the
global gradient. 
The benefit is that the gradients are exact and often very efficient to evaluate. 

\cite{scott2004response}, 
applied DDM to nonlinear beam-column elements.

\cite{al-aukaily2017response}




\subsection{Automatic Differentiation in Computational Mechanics}

In the past decade, automatic differentiation (AD) has made
significant inroads into computational mechanics, offering a way to
obtain exact gradients \emph{without} manual derivation. 
\Autodiff{} treats the finite element
algorithm as a composition of differentiable operations and applies the
chain rule systematically. 

Historic accounts of the development
of \autodiff{} can be found in \textcite{iri2009automatic}. 
%
Such treatments generally begin with the work of Wengert (e.g. \textcite{wengert1964simple}). 
A more recent work,
\textcite{elliott2009beautiful} develops an elegant presentation and
implementation of forward-mode AD in the context of purely functional
programming.

A thorough treatment of \autodiff{} for iterative procedures such as the use of
the Newton-Raphson method for solving nonlinear equations is presented
by \textcite{beck1994automatic}, \textcite{gilbert1992automatic} and
\textcite{griewank1993derivative}. 

Modern \autodiff{} tools can compute Jacobians and
gradients programmatically, and researchers have started
leveraging these tools in finite element contexts:

\begin{itemize}
\item
  \textbf{AD in FE Frameworks:} Large-scale multiphysics codes have
  begun to incorporate \autodiff{} to form exact system Jacobians. 
  A notable
  example is the MOOSE framework \citep{MOOSE} which uses the MetaPhysicL \autodiff{}
  library. 
  \cite{lindsay2021automatic} describe the integration of \autodiff{}
  into MOOSE's finite element routines for the automatic
  computation of consistent tangent matrices for nonlinear solves. 
  They report that this approach yields an exact Jacobian at relatively small overhead, enhancing the robustness and convergence of Newton--Raphson solvers. 
\item
  \textbf{Differentiable FE Libraries:} 
  \cite{liang2023pytorchfea} introduced
  \emph{PyTorch-FEA}, which implements finite element analysis using
  PyTorch tensor operations. 
  By ``taking advantage of autograd, an
  automatic differentiation mechanism in PyTorch,'' the library can
  seamlessly compute gradients of FE response with respect to input
  parameters. 
  This enabled them to tackle inverse problems by coupling
  the FE model with neural networks. 

  
  JAX-FEM \citep{xue2023jaxfem}  is a differentiable 3D FE solver built on Google's \Jax library \citep{frostig2018compiling}. 
  % It handled a 3D problem with 7.7 million
  % degrees of freedom and achieved an order-of-magnitude speedup over a
  % commercial FE code by running on GPU. 
  JAX-FEM uses reverse-mode \autodiff{} to
  obtain gradients, so one can perform automatic inverse design
  -- e.g.~topology optimization -- ``in a fully automatic manner
  without the need to manually derive sensitivities''. 
  % This
  % shows that with modern hardware and \autodiff{} tools, even large-scale FE
  % models can be made differentiable, opening the door to integrating FEM
  % with AI-driven design. 
  % (That said, JAX-FEM's success also relies on
  % problem structure and substantial computing power; fully
  % differentiating a million-DOF nonlinear problem is still extremely
  % memory-intensive in general.)
\item
  \textbf{AD in Specialized Methods:} Researchers have also applied \autodiff{}
  within specific computational mechanics techniques to simplify their
  implementation. 
  \cite{oberbichler2021efficient} used \autodiff{} to derive the internal force vector and consistent stiffness matrix for
  nonlinear isogeometric elements from a single energy
  functional. 
  They compared forward-mode vs.~reverse-mode
  \autodiff{} for this task and found the reverse (adjoint) mode far more
  efficient for high-degree elements: the reverse-mode Jacobian required
  significantly fewer operations and intermediate storage, and
  crucially, the number of AD intermediate variables no longer grew with
  the polynomial degree of shape functions. 
  In FE problems
  with many state variables but relatively fewer output quantities
  (e.g.~objective function or a few constraints), adjoint/reverse mode
  is usually optimal.
\end{itemize}

While \autodiff{} offers automation and exactness, it can
become computationally prohibitive for entire nonlinear finite element
simulations. 
The memory footprint of reverse-mode \autodiff{} is a known concern. 
In practice, fully
differentiable finite element codes often require careful engineering to be feasible. 
This is where hybrid approaches enter the picture -- combining analytic and
automatic differentiation to balance performance with flexibility.

\subsection{Hybrid Approaches}

Given the pros and cons of the two methods, a hybrid
differentiation strategy is attractive: use analytic
differentiation for parts of the problem where we can efficiently 
derive formulas, and use \emph{automatic differentiation} to handle the remainder. 
This idea has been explored in various forms.

% \textbf{Concept of Hybrid Differentiation:} 
\cite{bartlett2008hybrid} were early proponents of this approach in the context of
simulation codes. 
They observed that ``present AD tools\ldots rarely
compare to analytically derived versions'' in efficiency, but \autodiff{} is
far easier to implement. 
Thus, \emph{the focus of their work
was to combine the strength of these methods into a hybrid
strategy''}. 
In other words, one should ``select the
appropriate components'' of the simulation to differentiate manually
vs.~automatically to achieve an optimal balance of developer effort
and runtime cost. 
In their demonstration (fluid dynamics simulation in
C++), they hand-differentiated critical low-level routines and used
template-based AD for the rest -- yielding near-analytic performance
with much less manual coding.

% \textbf{Contemporary Examples:} 
Recently, \cite{dummer2024robust} put a hybrid approach into practice for nonlinear material modeling. 
They introduced a \emph{``semi-analytical split approach''}
  for a finite strain generalized continuum model. 
  In their framework, the global finite element equations were derived
  analytically as usual, but the material's stress-update
  algorithm was differentiated using automatic differentiation
  with hyper-dual numbers. 
  Hyper-dual \autodiff{} provided first and
  second derivatives of the stress update,
  yielding an exact consistent tangent matrix for the material
  sub-problem. 
  This split saved them from manually
  differentiating an arduous constitutive law, while still embedding
  that law into an efficient FE solver. 
  In effect, Dummer \emph{et al.} achieved a modular differentiation: the element
  stiffness contributions from a complicated material are obtained by 
  \autodiff{}, and everything else (element formulation, assembly, solver) uses
  standard analytic differentiation.

% \textbf{Adjoint-AD Hybrids:} 
Another approach to hybridization is mixing \emph{adjoint methods} with \autodiff{}. 
Adjoint (reverse) differentiation is itself an analytic approach at the system level, often used in topology optimization or sensitivity analysis to avoid recomputing a full Jacobian. 
One can use an adjoint equation to get the gradient of a single objective with respect to many inputs very efficiently -- but one still needs partial derivatives (Jacobian blocks) inside that adjoint equation. 
%
\cite{oberbichler2021efficient} apply the adjoint method to the global IGA
system, but they obtain the necessary
element-level derivatives via algorithmic differentiation of
the energy functional. 
The result was a clean separation: the
CAD-based geometry routines were automatically differentiated, and
then the adjoint method assembled the total sensitivity with far fewer
operations than a pure forward differentiation would require. This
hybrid of \autodiff{} and adjoint analytic methods allowed them to tackle shape
sensitivities for very high-order elements efficiently. 


% \subsection{Objectives}
% We aim to differentiate the FE element kernel
% analytically for efficiency, and let a higher-level \autodiff{} system handle
% glue code and complex couplings for flexibility.

% Pure AD-based FEM implementations exist and have shown what's possible (e.g.~end-to-end differentiable simulators like JAX-FEM and PyTorch-FEA). 
% However, they also highlight the downsides -- namely, the runtime and memory costs -- when applied to large, nonlinear FE problems. 
% On the other hand, decades of work on DDM in computational mechanics have furnished methods to get
% gradients efficiently for very specific problems (e.g.~the OpenSees DDM
% for structural sensitivities, or recent DDM for geometrically-nonlinear
% frame elements). 

% The hybrid approach seeks to get the ``best of
% both worlds.'' 
% By embedding an analytically differentiated finite element
% kernel inside a higher-level \autodiff{} framework, we retain the
% efficiency of hand-coded derivatives for the heavy lifting (the finite
% element internal calculations) and the flexibility of automatic
% differentiation for coupling with other software (optimization
% algorithms, machine learning models, etc.). 
% This approach has only started to appear in the literature in the last few years, and to our knowledge, no general-purpose Python-based finite element tool with built-in DDM
% sensitivities exists yet. 
% Our work will fill that gap. 
% We will demonstrate how a DDM-enhanced finite element solver can be embedded
% in an \autodiff{} environment to enable gradient-based design,
% optimization, and learning on complex FE models -- without the
% prohibitive memory cost of differentiating through the entire nonlinear
% solve. 
% This aligns with the trends seen in recent research and pushes
% them a step further by uniting the efficiency of classical methods with
% the flexibility of modern \autodiff{} frameworks.






\section{Setting. State and response maps}


% Let \(\ModelFamily\) denote a fixed discrete model family. 
% It determines the governing residual and the output operator, but it is not itself the differentiated map. 


A system that relies entirely on direct differentiation must implement differentiation rules for any parameterization of the model that may be of interest. 
In the proposed approach, direct differentiation is only required for a minimal parameter set that is required to perform the analysis. 
These parameters are referred to as \emph{core parameters} $\corey \in \CoreSpace\sim\mathbb{R}^{\ncore}$.  
Typical core parameters include material properties, fiber coordinates and areas, nodal
coordinates, and prescribed force magnitudes. 


Any parameter set from which the core parameter may be derived is referred to as the \emph{host parameterization} $\haty \in \HostSpace\sim\mathbb{R}^{\nhost}$. 
These generally represent quantities that the analyst wishes to differentiate with respect to and typically include  cross-section dimensions, load
multipliers, random-field coefficients, or any other inputs that the
user controls.


A user-defined parameterization is a differentiable map
% \[
% \ParameterMap : \HostSpace \to \CoreSpace,
% \qquad
% \CoreParameters = \ParameterMap(\HostParameters).
% \]
\begin{equation}\label{eq:Ymap}
    \Ymap : \mathbb{R}^{\nhost} \to \mathbb{R}^{\ncore},
    \qquad \haty \mapsto \corey = \Ymap(\haty),
\end{equation}
For each \(\CoreParameters \in \CoreSpace\), the equilibrium state \(\State \in \StateSpace\) satisfies
\[
\Residual_{\ModelFamily}(\State,\CoreParameters) = \bm 0.
\]
where \(\ModelFamily\) is a fixed discretization of the physical model. 
For a well posed problem this implicitly defines a differentiable solution map:
\[
\StateMap_{\ModelFamily} : \CoreSpace \to \StateSpace,
\qquad
\State = \StateMap_{\ModelFamily}(\CoreParameters).
\]
Given a quantity of interest \(
\Observation_{\ModelFamily} : \StateSpace \times \CoreSpace \to \ResponseSpace
\) 
the reduced response map is
\begin{equation}
\ResponseMap_{\ModelFamily}(\CoreParameters)
=
\Observation_{\ModelFamily}(\StateMap_{\ModelFamily}(\CoreParameters),\CoreParameters).
\label{eq:def-reduced-map}
\end{equation}
The user-facing response is the composite map
\[
\ComposedResponse_{\ModelFamily}
=
\ResponseMap_{\ModelFamily}\circ\ParameterMap
:
\HostSpace \to \ResponseSpace,
\qquad
\CoreParameters=\ParameterMap(\HostParameters).
\]
\begin{remark}
If the quantity of interest is the full discrete state, then one may take
\[
\Observation_{\ModelFamily}(\State,\CoreParameters)=\State,
\]
in which case
\[
\ResponseMap_{\ModelFamily}=\StateMap_{\ModelFamily}.
\]
\end{remark}

\section{Differential Structure}

For any differentiable map \(f:X\to Z\), define its tangent extension by
\[
\Tan{f}(x,\dot x)
=
\bigl(
f(x),
\D f(x)[\dot x]
\bigr).
\]
% This is the natural primal-tangent object in forward mode. 
% It returns both the value and the Jacobian-vector product.
For composable differentiable maps \(g:X\to Y\) and \(f:Y\to Z\),
\[
\Tan{f\circ g} = \Tan{f}\circ\Tan{g}.
\]

The total derivative of the reduced map \eqref{eq:def-reduced-map} satisfies the chain rule
\begin{equation}
\label{eq:composed-derivative}
\D \ComposedResponse_{\ModelFamily}(\HostParameters)
=
\D \ResponseMap_{\ModelFamily}(\CoreParameters)
\circ
\D \ParameterMap(\HostParameters)
\in \mathcal L(\HostSpace,\ResponseSpace).
\end{equation}
After choosing coordinates on \(\HostSpace\), \(\CoreSpace\), and \(\ResponseSpace\), this becomes
\begin{equation}
\label{eq:composed-jacobian}
\frac{\partial \ComposedResponse_{\ModelFamily}}{\partial \HostParameters}
=
\frac{\partial \ResponseMap_{\ModelFamily}}{\partial \CoreParameters}
\;
\frac{\partial \ParameterMap}{\partial \HostParameters}.
\end{equation}

\paragraph{Kernel level operators.}

The finite element kernel is required to provide the tangent extension of the reduced response map
\[
\Tan{\ResponseMap_{\ModelFamily}} :
\CoreSpace \times \CoreSpace
\to
\ResponseSpace \times \ResponseSpace.
\]
% Equivalently, for any \((\CoreParameters,\dot{\CoreParameters})\), the kernel returns
% \[
% \ResponseMap_{\ModelFamily}(\CoreParameters),
% \qquad
% \D \ResponseMap_{\ModelFamily}(\CoreParameters)[\dot{\CoreParameters}].
% \]

If \(\State=\StateMap_{\ModelFamily}(\CoreParameters)\), the tangent state \(\dot{\State}\) is defined by the linearized equilibrium equation
\[
\D_1 \Residual_{\ModelFamily}(\State,\CoreParameters)[\dot{\State}]
=
-
\D_2 \Residual_{\ModelFamily}(\State,\CoreParameters)[\dot{\CoreParameters}].
\]
The corresponding response variation is
\[
\D \ResponseMap_{\ModelFamily}(\CoreParameters)[\dot{\CoreParameters}]
=
\D_1 \Observation_{\ModelFamily}(\State,\CoreParameters)[\dot{\State}]
+
\D_2 \Observation_{\ModelFamily}(\State,\CoreParameters)[\dot{\CoreParameters}].
\]
When the quantity of interest is a pure state extraction, the second term vanishes.

\paragraph{Host level operators.}

The host system is required to provide the tangent extension of the parameterization map
\[
\Tan{\ParameterMap} :
\HostSpace \times \HostSpace
\to
\CoreSpace \times \CoreSpace.
\]
Equivalently, for any \((\HostParameters,\dot{\HostParameters})\), the host returns
\[
\ParameterMap(\HostParameters),
\qquad
\D \ParameterMap(\HostParameters)[\dot{\HostParameters}].
\]

This layer contains all differentiable preprocessing that is external to the finite element kernel. It includes geometry generation, section construction, discretization rules, constitutive preprocessing, and load parametrization.

% \section{Differential structure}
% \label{sec:differential}

% The total derivative of the reduced map \eqref{eq:def-reduced-map} satisfies the chain rule
% \begin{equation}
% \label{eq:composed-derivative}
% \D \ComposedResponse_{\ModelFamily}(\HostParameters)
% =
% \D \ResponseMap_{\ModelFamily}(\CoreParameters)
% \circ
% \D \ParameterMap(\HostParameters)
% \in \mathcal L(\HostSpace,\ResponseSpace).
% \end{equation}
% After choosing coordinates on \(\HostSpace\), \(\CoreSpace\), and \(\ResponseSpace\), this becomes
% \begin{equation}
% \label{eq:composed-jacobian}
% \frac{\partial \ComposedResponse_{\ModelFamily}}{\partial \HostParameters}
% =
% \frac{\partial \ResponseMap_{\ModelFamily}}{\partial \CoreParameters}
% \;
% \frac{\partial \ParameterMap}{\partial \HostParameters}.
% \end{equation}
The Jacobian in \eqref{eq:composed-jacobian} is not formed explicitly. One evaluates its action on tangent vectors or on cotangent vectors. The forward view is encoded by
\[
\Tan{\ComposedResponse_{\ModelFamily}}
=
\Tan{\ResponseMap_{\ModelFamily}}
\circ
\Tan{\ParameterMap},
\]
and the reverse view is encoded by
\[
\D \ComposedResponse_{\ModelFamily}(\HostParameters)^*
=
\D \ParameterMap(\HostParameters)^*
\circ
\D \ResponseMap_{\ModelFamily}(\CoreParameters)^*.
\]

\subsection{Forward mode}
\label{sec:jvp}

Given a tangent direction \(\dot{\HostParameters}\in \HostSpace\), the Jacobian--vector product is
\begin{equation}
\label{eq:jvp-composed}
\dot{\bm z}
=
\D \ComposedResponse_{\ModelFamily}(\HostParameters)[\dot{\HostParameters}]
=
\underbrace{\D \ResponseMap_{\ModelFamily}(\CoreParameters)}_{\text{kernel}}
\Bigl[
\underbrace{\D \ParameterMap(\HostParameters)[\dot{\HostParameters}]}_{\text{host}}
\Bigr],
\end{equation}
where \(\bm z \coloneq \ComposedResponse_{\ModelFamily}(\HostParameters)\). 
The computation is performed into two stages:
\begin{enumerate}
    \item \textbf{Host forward pass.} Evaluate
    \[
    (\CoreParameters,\dot{\CoreParameters})
    =
    \Tan{\ParameterMap}(\HostParameters,\dot{\HostParameters}),
    \qquad
    \dot{\CoreParameters}
    =
    \D \ParameterMap(\HostParameters)[\dot{\HostParameters}].
    \]

    \item \textbf{Kernel forward pass.} Evaluate
    \[
    (\bm z,\dot{\bm z})
    =
    \Tan{\ResponseMap_{\ModelFamily}}(\CoreParameters,\dot{\CoreParameters}),
    \qquad
    \bm z=\ResponseMap_{\ModelFamily}(\CoreParameters).
    \]
\end{enumerate}

When the response map is induced by the residual and observation operators,
\[
\Residual_{\ModelFamily}(\State,\CoreParameters)=\bm 0,
\qquad
\bm z=\Observation_{\ModelFamily}(\State,\CoreParameters),
\]
the tangent variables satisfy
\begin{align}
\label{eq:forward-linearized-equilibrium}
\D_1 \Residual_{\ModelFamily}(\State,\CoreParameters)[\dot{\State}]
&=
-
\D_2 \Residual_{\ModelFamily}(\State,\CoreParameters)[\dot{\CoreParameters}],
\\
\label{eq:forward-response-variation}
\dot{\bm z}
&=
\D_1 \Observation_{\ModelFamily}(\State,\CoreParameters)[\dot{\State}]
+
\D_2 \Observation_{\ModelFamily}(\State,\CoreParameters)[\dot{\CoreParameters}].
\end{align}
After choosing coordinates \(\CoreParameters=(y_1,\ldots,y_{n_{\mathrm c}})\),
\begin{equation}
\label{eq:jvp-contraction}
\dot{\bm z}
=
\sum_{i=1}^{n_{\mathrm c}}
\frac{\partial \ResponseMap_{\ModelFamily}}{\partial y_i}(\CoreParameters)\,\dot y_i.
\end{equation}
Forward mode is natural when derivatives are required in only a small number of host directions.

\subsection{Reverse mode}
\label{sec:vjp}

Given a cotangent \(\bar{\bm z}\in \ResponseSpace^*\), the vector--Jacobian product is
\begin{equation}
\label{eq:vjp-composed}
\bar{\HostParameters}
=
\D \ComposedResponse_{\ModelFamily}(\HostParameters)^*[\bar{\bm z}]
=
\underbrace{\D \ParameterMap(\HostParameters)^*}_{\text{host}}
\Bigl[
\underbrace{\D \ResponseMap_{\ModelFamily}(\CoreParameters)^*[\bar{\bm z}]}_{\text{kernel}}
\Bigr].
\end{equation}
In finite-dimensional Euclidean coordinates this is
\begin{equation}
\label{eq:vjp-euclidean}
\bar{\HostParameters}
=
\bigl[\D \ParameterMap(\HostParameters)\bigr]^{\top}\bar{\CoreParameters},
\qquad
\bar{\CoreParameters}
=
\bigl[\D \ResponseMap_{\ModelFamily}(\CoreParameters)\bigr]^{\top}\bar{\bm z}.
\end{equation}
For a scalar objective \(\Objective:\ResponseSpace\to\mathbb R\), one takes
\[
\bar{\bm z}
=
\D \Objective(\bm z),
\qquad
\bar{\HostParameters}
=
\D (\Objective\circ\ComposedResponse_{\ModelFamily})(\HostParameters),
\]
so that a single reverse sweep delivers the gradient with respect to all host parameters.

The computation again separates into two stages:
\begin{enumerate}
    \item \textbf{Kernel reverse pass.} Compute the core-parameter cotangent
    \[
    \bar{\CoreParameters}
    =
    \D \ResponseMap_{\ModelFamily}(\CoreParameters)^*[\bar{\bm z}].
    \]
    After choosing coordinates \(\CoreParameters=(y_1,\ldots,y_{n_{\mathrm c}})\),
    \begin{equation}
    \label{eq:vjp-kernel-contraction}
    \bar y_i
    =
    \left(
        \frac{\partial \ResponseMap_{\ModelFamily}}{\partial y_i}(\CoreParameters)
    \right)^{\!\top}
    \bar{\bm z},
    \qquad
    i=1,\ldots,n_{\mathrm c}.
    \end{equation}
    If \(\ResponseMap_{\ModelFamily}\) is described through
    \(\Residual_{\ModelFamily}\) and \(\Observation_{\ModelFamily}\), the same quantity may be obtained from the adjoint system
    \begin{align}
    \label{eq:adjoint-equation}
    \D_1 \Residual_{\ModelFamily}(\State,\CoreParameters)^*[\AdjointState]
    &=
    \D_1 \Observation_{\ModelFamily}(\State,\CoreParameters)^*[\bar{\bm z}],
    \\
    \label{eq:adjoint-gradient}
    \bar{\CoreParameters}
    &=
    \D_2 \Observation_{\ModelFamily}(\State,\CoreParameters)^*[\bar{\bm z}]
    -
    \D_2 \Residual_{\ModelFamily}(\State,\CoreParameters)^*[\AdjointState].
    \end{align}

    \item \textbf{Host reverse pass.} Propagate \(\bar{\CoreParameters}\) through the parameter map:
    \[
    \bar{\HostParameters}
    =
    \D \ParameterMap(\HostParameters)^*[\bar{\CoreParameters}].
    \]
\end{enumerate}
Reverse mode is natural when the number of outputs, or the number of scalar objectives, is small relative to the number of host parameters.

\begin{remark}
If the kernel derivative is available in coordinates in the form 
\(\{\partial \ResponseMap_{\ModelFamily}/\partial y_i\}_{i=1}^{n_{\mathrm c}}\),
then both \eqref{eq:jvp-contraction} and \eqref{eq:vjp-kernel-contraction}
are contractions against the same family. The principal distinction between forward and reverse modes therefore arises in the differentiation of the host map \(\ParameterMap\).
\end{remark}


% \section{Summary and Implementation}
% % \section{Composition Protocol}\label{sec:protocol}

% The response exposed to the user is the composite map
% \[
% \ComposedResponse_{\ModelFamily}
% =
% \ResponseMap_{\ModelFamily}\circ\ParameterMap
% :
% \HostSpace \to \ResponseSpace.
% \]
% Its tangent extension follows from the compositional identity,
% \[
% \Tan{\ComposedResponse_{\ModelFamily}}
% =
% \Tan{\ResponseMap_{\ModelFamily}}\circ\Tan{\ParameterMap}.
% \]
% In directional derivative form,
% \[
% \D \ComposedResponse_{\ModelFamily}(\HostParameters)[\dot{\HostParameters}]
% =
% \D \ResponseMap_{\ModelFamily}(\ParameterMap(\HostParameters))
% \Bigl[
% \D \ParameterMap(\HostParameters)[\dot{\HostParameters}]
% \Bigr].
% \]

% This identity is the central statement of the formulation. The host differentiates the construction of native parameters. The kernel differentiates the mechanics response with respect to those parameters. The full derivative is obtained by composition, and neither layer needs to materialize a full Jacobian.

% % \section{Scalar objectives}

% If a scalar objective is introduced,
% \[
% \Objective : \ResponseSpace \to \mathbb R,
% \qquad
% \widehat{\Objective}
% =
% \Objective \circ \ComposedResponse_{\ModelFamily},
% \]
% then forward differentiation gives
% \[
% \D \widehat{\Objective}(\HostParameters)[\dot{\HostParameters}]
% =
% \D \Objective\bigl(\ComposedResponse_{\ModelFamily}(\HostParameters)\bigr)
% \Bigl[
% \D \ComposedResponse_{\ModelFamily}(\HostParameters)[\dot{\HostParameters}]
% \Bigr].
% \]
